%% kapa: exam=1 round=20 year=2009 period=1 form=A subject=economics subject_ko=경제원론 seq=36 number=76
%% source: Q-Net 감정평가사 기출문제 (cst00309) artlSeq=5024178 file=extracted/1008/5024178/2009년 제20회 감정평가사 1차시험 기출문제/1교시.hwp
%% conversion: kordoc 4.24.51 -> tools/gen1.py (auto); review=proofread-sonnet; flags=none
\begin{question}{76}
인플레이션과 관련된 설명으로 옳지 \uline{않은} 것은?
\begin{choices}
  \item 폐쇄경제에서 완전고용상태일 때 총수요가 총공급을 초과하면 인플레이션이 발생한다.
  \item 인플레이션이 예상될 때 개인들이 재화를 미리 사두려고 하면 인플레이션은 더욱 심화된다.
  \item 수입원자재 가격의 상승은 비용인상 인플레이션을 유발할 수 있다.
  \item 인플레이션으로 사회구성원 사이에 소득이나 부(富)가 재분배되기도 한다.
  \item 중앙은행은 인플레이션을 진정시키기 위해 국공채를 매입한다.
\end{choices}
\end{question}
\answer{5}
